% !TEX root = ../auv-thesis.tex
% Authoritative thesis source; edit this TeX file directly.
% Legacy Markdown material is reference-only and must not overwrite this file.

\chapter{声磁协同状态估计}

\section{引言}\label{sec:ch03-3-1}

\subsection{水下复杂环境下感知面临的挑战}\label{sec:ch03-3-1-1}

在讨论具体算法之前，应先把水下状态估计这件事的难度说清楚。陆上机器人之所以能凭借 GNSS、轮速里程计和稳定的视觉特征获得近乎``开盒即用''的定位结果，本质上得益于三个先决条件：全局参考可用、里程计直接可测、视觉环境几何稳定。把同样的需求搬到水下，会发现这三个条件全部不再成立，GNSS 信号被海水迅速衰减、轮速等机械里程计无从谈起、可见光成像在浑浊水中很快失去用途。这意味着水下状态估计从一开始就只能在``有限可用、间歇可用''的传感器组合上构建。

更进一步看，水下电缆巡检对状态估计的要求并不只是``知道自己在哪里''。它还要回答两个并列的问题:一是 AUV 相对于电缆的几何关系如何（例如横向偏移、相对航向、离底高度），二是当前对这些估计有多大把握（也就是不确定性度量）。第一类问题决定控制器的参考量是否合理，第二类问题决定行为树是否要切换为更保守的模式。如果只解决第一类问题而忽略第二类，控制器就会把不可靠的估计当作真值，使整个系统变得脆弱。

把这两点叠加到真实海况上，问题进一步具体化：DVL 在浊水或海草遮挡下会出现长达数秒的``底部跟踪丢失''；声呐在底质突变或多径反射明显的区域会输出大片杂波，使电缆边缘特征被淹没；磁场观测在三相抵消严重的方位、铠装屏蔽集中的位置或外部电磁干扰下会出现畸变；通信链路在母船-AUV 之间的可用带宽本就有限，进一步限制了可以传回的中间数据量。这些情况并不是``偶发故障''，而是水下作业的常态。

由此可以得出一个对后续设计具有决定性的认识：状态估计器不应假设传感器始终可靠，而应在传感器条件变化时维持状态估计连续性，并把不确定性显式传递给控制和决策层。本章的所有讨论，从多源数据预处理、融合算法设计，到不确定性量化与置信度衰减，都建立在这一前提之上。

\subsection{组合导航系统总体框架设计}\label{sec:ch03-3-1-2}

延续上节``传感器随时可能不靠谱''的认识，本文把组合导航系统设计为``惯性递推 + DVL/深度/声磁观测修正 + 不确定性在线评估''三者并行的结构，而不是单纯的``传感器加权融合''。其中，以惯性递推为高频骨架、以 DVL 等低频观测抑制漂移的基本思路建立在水下组合导航研究基础上\cite{hegrenaes2009doppler,zhang2023sinsdvl}；本文进一步把声磁观测质量及其不确定性传递纳入统一结构。三者之间的耦合方式如下：递推保证状态在没有观测时仍然连续；观测在到达时被加权修正；不确定性评估则决定每次观测的修正幅度。

本文选择误差状态扩展卡尔曼滤波（error-state extended Kalman filter, ES-EKF）作为融合骨架，而不是直接采用标准扩展卡尔曼滤波（extended Kalman filter, EKF）。这一选择的理论基础来自四元数误差状态建模及其在惯性组合导航中的应用\cite{sola2017quaternion,madyastha2011eskf}，本文则针对水下任务作出具体取舍。第一，AUV 在巡检过程中会经历较大幅度的航向转动（例如之字形扫描的折返、绕过障碍的转弯），因而名义姿态采用单位四元数表示以避免欧拉角奇异；ES-EKF 进一步用局部小角度扰动表示姿态误差，使线性化始终发生在小量附近，并便于维持四元数单位约束。第二，ES-EKF 在``误差状态接近零''的工作点附近做线性化，适合水下高频 IMU 递推与低频观测修正并存的场景。本节先说明算法选择及其输入输出关系，名义状态、误差状态和观测更新的形式化定义统一放在 3.3.1 节给出。

整个组合导航系统的输入输出可以概括为：输入端接收 IMU、DVL、深度、声呐特征、磁场反演结果以及它们各自的质量度量；输出端给出位置、速度、姿态、传感器偏置、对应协方差矩阵，以及由协方差派生的标量置信度。这个置信度不仅用于内部自适应观测加权，还会被传递给行为树和 UA-MPC 控制层，让上层决策与控制能够``知道当前估计有多大把握''。换句话说，本系统的设计哲学是把不确定性视作和状态本身同等重要的输出，而不是只在内部使用的中间量。该共享信号从物理扰动到安全执行的传递链路如图~\ref{fig:ch03-uncertainty-highway}所示。

\begin{figure}[htbp]
\centering
\includegraphics[width=0.95\linewidth,height=\textheight,keepaspectratio,alt={感知不确定性从物理扰动到安全执行的传递链路（状态估计、行为树与 UA-MPC 的共享置信度接口）}]{architecture/auv_v2_uncertainty_highway.png}
\caption{感知不确定性从物理扰动到安全执行的传递链路（状态估计、行为树与 UA-MPC 的共享置信度接口）}\label{fig:ch03-uncertainty-highway}
\end{figure}
\FloatBarrier

\section{多源异构数据预处理与时空同步}\label{sec:ch03-3-2}

\subsection{基于 ROS2 的异步传感器时间戳对齐算法}\label{sec:ch03-3-2-1}

把状态估计的整体框架确定下来之后，第一件需要解决的工程问题就是``不同传感器以不同频率到达，如何让它们在同一个时间轴上参与计算''。这个问题的特殊之处在于：水下传感器的频率差异不是几倍，而是一个数量级以上。本系统中，TMR8637 选型表给出的响应带宽为 DC--300 kHz，但当前 SK2301 采集配置以名义 2000 Hz 完成三轴原始采集，并以 200 个样本为步长更新数字信号处理窗口；IMU 约 100 Hz、深度约 50 Hz、DVL 仅约 2 Hz。这里必须区分器件带宽、ADC 名义采样率、数字特征更新率和融合更新率，四者不能用同一个``采样频率''替代。当前真机链路的有效采样率仍存在低于名义值的诊断项，因而时间轴标定属于实物部署前必须闭合的接口问题。如果状态估计器简单地把``最近一次到达的观测''当作当前观测使用，就会产生两类错误：一是把已经过期数百毫秒的 DVL 速度当作当前速度；二是把不同时刻的观测错位组合，从而引入虚假相关性。

更进一步，水下任务中观测延迟的来源不止采样频率本身。AMD/VxWorks 经 UDP 上行的传感器帧带有自身的网络时延和抖动，仿真后端在调度时也存在批处理延迟。这意味着``采样时间''和``到达时间''并不重合，必须显式区分。本系统在 Mock AMD 子系统中通过两个解耦模块来描述这件事：\texttt{Transport\allowbreak{}Delay\allowbreak{}Queue} 模拟从采样到达消费端的非确定性传输延迟（带均值和抖动两个参数），\texttt{Sensor\allowbreak{}Sample\allowbreak{}Cache} 则在采样侧维护各传感器独立的``最后采样时间''和缓存值。

\texttt{Sensor\allowbreak{}Sample\allowbreak{}Cache} 的工作方式很直白：对每个传感器维护一个采样周期，每帧调用时检查是否超过该周期，满足则更新缓存值并打上新的时间戳，否则继续返回上次缓存（hold）。DVL 额外加入``底部跟踪丢失''开关，在丢失窗口内即使到了采样周期也保持上次值不变，模拟真实 DVL 在浊水或杂草中的失效行为。这种实现方式有意把``频率差''和``丢失''两件事用同一套接口表达，使得 ES-EKF 上层无需关心底层失效模式，只需根据时间戳和有效性标志做对齐和加权。

为把上述``频率差 + 传输延迟''落到可量化的工程指标，先归纳本系统各传感器的采样频率与传输时延特征：

\begin{longtable}[]{@{}
  >{\raggedright\arraybackslash}p{(\linewidth - 6\tabcolsep) * \real{0.2500}}
  >{\raggedright\arraybackslash}p{(\linewidth - 6\tabcolsep) * \real{0.2500}}
  >{\raggedright\arraybackslash}p{(\linewidth - 6\tabcolsep) * \real{0.2500}}
  >{\raggedright\arraybackslash}p{(\linewidth - 6\tabcolsep) * \real{0.2500}}@{}}
\caption{传感器系统更新率与传输时延特征}\label{tab:ch03-sensor-rates}\tabularnewline
\toprule\noalign{}
\begin{minipage}[b]{\linewidth}\raggedright
传感器
\end{minipage} & \begin{minipage}[b]{\linewidth}\raggedright
采样频率
\end{minipage} & \begin{minipage}[b]{\linewidth}\raggedright
典型传输时延
\end{minipage} & \begin{minipage}[b]{\linewidth}\raggedright
链路/接口
\end{minipage} \\
\midrule\noalign{}
\endfirsthead
\toprule\noalign{}
\begin{minipage}[b]{\linewidth}\raggedright
传感器
\end{minipage} & \begin{minipage}[b]{\linewidth}\raggedright
采样频率
\end{minipage} & \begin{minipage}[b]{\linewidth}\raggedright
典型传输时延
\end{minipage} & \begin{minipage}[b]{\linewidth}\raggedright
链路/接口
\end{minipage} \\
\midrule\noalign{}
\endhead
\bottomrule\noalign{}
\endlastfoot
TMR8637 三轴 TMR 磁传感器 & 2000 Hz 名义原始采集；10 Hz 特征步进 & 待按有效采样率与 TCP 链路实测 & SK2301 24-bit ADC / Ethernet TCP / ROS2 \\
IMU & 100 Hz & \(< 2\text{ ms}\) & RS232 \\
DVL & 2 Hz & \(\approx 50\text{ ms}\) & RS422 / UDP \\
侧扫/浅剖 EdgeTech 2205 & 2--5 Hz & \(150\sim300\text{ ms}\) & 1000Base-T / TCP \\
\end{longtable}

时延为什么必须显式补偿？做一个量级估算即可看清：当 AUV 以 \(v = 1.0\text{ m/s}\) 巡航、某观测延迟 \(\tau_{delay} = 200\text{ ms}\) 时，若直接把它当作当前时刻观测使用，引入的空间错位为

\[\Delta x_{error} = v \cdot \tau_{delay} = 1.0\text{ m/s} \times 0.2\text{ s} = 0.2\text{ m}\]

恰好触及 T/CSGPC 规定的埋深/路由误差上限。本系统据此采用``硬件锁相 + 软件回溯''的双重同步：其一是硬件层面的 1PPS 秒脉冲锁相，由 DVL/GPS 产生的 1 Hz 方波接入声呐 \texttt{Trig1} 与 Jetson GPIO，各节点在上升沿清零内部计数器，消除晶振漂移（对应第 2.5.1 节的物理基础）；其二是软件层面为每帧观测打上 ROS2 Header 的 Unix 纳秒时间戳，并按硬件时延修正回采样时刻：

\[t_{stamp} = t_{system\_now} - \tau_{hardware\_latency}\]

在此基础上处理``低频高延迟观测滞后到达''的问题：为高频状态开辟一个保存过去 \(2.0\text{ s}\)（容量 \(N = 400\)）的双向环形缓冲区；当声呐/视觉等低频观测在 \(t_{recv}\) 到达时，先提取其真实发射时间 \(t_{ping} = t_{recv} - \tau_{delay}\)，在缓冲区中检索满足 \(t_k \le t_{ping} < t_{k+1}\) 的相邻两帧，做一阶线性（或 Hermite）插值得到发射时刻的状态：

\[\hat{\boldsymbol{x}}(t_{ping}) = \boldsymbol{x}_k + \frac{t_{ping} - t_k}{t_{k+1} - t_k}\left(\boldsymbol{x}_{k+1} - \boldsymbol{x}_k\right)\]

随后在该历史位姿处执行观测更新步，再用 IMU 增量把状态推回当前时刻。经此回溯-插值处理，多源时空同步误差可控制在 \(< 1\text{ ms}\)（对应空间误差 \(< 1\text{ mm}\)），从根本上消除了因异步到达导致的虚假相关性。

需要诚实说明的是，仿真期通过软件模拟的``采样时间戳 + 缓存接口''已经能够支撑算法层面的多速率融合验证，但 PTP 硬件时戳、AMD 与 Jetson 之间真实的时钟漂移补偿等问题仍需要在实物部署阶段进一步完成。本文在写作时把这部分明确划入``接口已识别、硬件级证据待补''的边界，避免把仿真结果直接外推为实物结论。

\subsection{传感器安装偏差角标定与补偿}\label{sec:ch03-3-2-2}

时间对齐之后，下一个需要处理的是空间对齐。空间一致性问题的本质是：不同传感器物理上安装在 AUV 的不同位置，并且各自有自己的坐标系朝向，但是状态估计器只在一个统一的``AUV 参考坐标系''内做计算。如果不做坐标变换，两件事会同时出错，观测的几何位置错了（杆臂偏移），观测的方向也错了（坐标轴未对齐）。两者叠加后产生的系统性误差通常会被滤波器误判为``过程噪声大''，进而拉大协方差，整个估计精度下降。

为什么对水下电缆巡检任务这一问题尤其重要？关键在于磁传感器的安装方式。为了避开 AUV 本体钢结构的剩磁干扰，磁传感器通常装在 AUV 前端伸出的非磁性长杆末端，距离重心 0.5--1 m 量级。这种安装方式带来两个连带效应：第一，磁传感器实际测量点不在 AUV 重心，必须做杆臂改正（lever-arm correction）；第二，长杆受水流冲击会产生微小弯曲和扭转，使得安装方向标定不再是一次性工作，而需要周期性复核。

\begin{longtable}[]{@{}
  >{\raggedright\arraybackslash}p{(\linewidth - 6\tabcolsep) * \real{0.2500}}
  >{\raggedright\arraybackslash}p{(\linewidth - 6\tabcolsep) * \real{0.2500}}
  >{\raggedright\arraybackslash}p{(\linewidth - 6\tabcolsep) * \real{0.2500}}
  >{\raggedright\arraybackslash}p{(\linewidth - 6\tabcolsep) * \real{0.2500}}@{}}
\caption{传感器安装偏差与补偿状态}\label{tab:ch03-alignment}\tabularnewline
\toprule\noalign{}
\begin{minipage}[b]{\linewidth}\raggedright
偏差类型
\end{minipage} & \begin{minipage}[b]{\linewidth}\raggedright
影响
\end{minipage} & \begin{minipage}[b]{\linewidth}\raggedright
补偿方式
\end{minipage} & \begin{minipage}[b]{\linewidth}\raggedright
当前状态
\end{minipage} \\
\midrule\noalign{}
\endfirsthead
\toprule\noalign{}
\begin{minipage}[b]{\linewidth}\raggedright
偏差类型
\end{minipage} & \begin{minipage}[b]{\linewidth}\raggedright
影响
\end{minipage} & \begin{minipage}[b]{\linewidth}\raggedright
补偿方式
\end{minipage} & \begin{minipage}[b]{\linewidth}\raggedright
当前状态
\end{minipage} \\
\midrule\noalign{}
\endhead
\bottomrule\noalign{}
\endlastfoot
安装位置偏差 & 磁场观测点偏移 AUV 重心 & 杆臂改正（位移向量补偿） & 待实验 \\
安装偏差角 & 传感器坐标系与 AUV 坐标系未对齐 & 旋转矩阵补偿 & 待实验 \\
DVL 安装角 & DVL 波束方向与 AUV 体轴偏差 & 速度坐标变换 & 架构已识别 \\
\end{longtable}

DVL 安装角同样不可忽视。DVL 输出的是相对底部的速度向量，但其内部坐标系与 AUV 体坐标系有一个固定的旋转偏差（通常出厂时已经标定，但每次拆装都需要重新确认）。这个偏差如果不补偿，沿前进方向的速度估计会泄漏到横向，使滤波器看到一个``莫名其妙的横向速度''，进而干扰位置积分。

\subsection{磁场信号的工频提取与噪声抑制}\label{sec:ch03-3-2-3}

时空对齐解决了``对得上''的问题，下一个层次的问题则是``信号本身能不能用''。这一节的核心矛盾在于：海缆磁探测的有效信号可能本来就很弱（三相抵消 + 铠装屏蔽 + 螺旋绞合三重削弱），同时还会被地磁缓变、宽带电子噪声、AUV 自身磁噪声以及与目标同频但不同相干来源的工频串扰污染。如果不做预处理直接进入观测模型，信噪比会低到使观测协方差膨胀失控。

信号处理思路上有三层。第一层以待测海缆电流的参考频率 \(f_0\) 为中心做前置带宽限制，抑制低频地磁起伏和高频宽带噪声；第二层通过相敏正交解调提取与参考同频、同相干的目标分量，不能把目标工频本身作为陷波对象；第三层利用 AUV 轨迹与电缆空间响应的相关性进一步排斥非相干同频串扰。实际电网工频通常为 50 Hz 或 60 Hz，本文可控电缆和注入实验采用 45 Hz，以避开环境工频并保持参考源可控。若无法取得电缆电流参考，相干性与目标频率漂移则应作为观测质量的一部分进入不确定性评估。

但是，仅写``加滤波器''还不够。在状态估计的语境下，预处理的输出不只是干净信号本身，还应包括``信号有多干净''这一附加信息。具体而言，滤波器输出需要伴随一个估计的瞬时信噪比（SNR）或残差能量比，这个量进一步喂给后续的观测协方差自适应模块，SNR 高时观测协方差小、滤波器更信任磁观测，SNR 低时观测协方差大、滤波器更倾向于依赖预测。这种设计思想让``预处理''和``融合''从两个独立步骤变成一个连续链条。

模组层已有的物理输入证据来自编号 6\# TMR8637 随附三轴测试报告。该报告给出 X、Y、Z 三轴在 1 Hz 处的磁场噪声谱密度均为 200 pT/\(\sqrt{\text{Hz}}\)，其对数谱曲线在目标频率 50 Hz 附近约为 20--30 pT/\(\sqrt{\text{Hz}}\)，详见图~\ref{fig:ch02-tmr8637-noise-psd}。这一结果可作为 DLIA 输入端传感器本体噪声的先验量级，但还不是 TMR8637、SK2301、隔离供电与数字处理组成的完整采集链噪声；实际 SNR 仍应由运行时窗口内的同频信号和邻近频带噪声共同估计。

\textbf{从宽带滤波到相敏数字锁相放大。} 传统 FIR/IIR 带通滤波器受品质因数 \(Q = f_0/\Delta f\) 的限制，要把等效噪声带宽压到 1 Hz 以下通常需要较高阶数，并引入较长相位延迟。数字锁相放大的正交解调原理及数字实现可参见相关研究\cite{song2020dlia}。基于这一通用方法，本系统采用数字锁相放大器（digital lock-in amplifier, DLIA），在通用推导中保留参考频率 \(f_0\)，并在本文实测中取 \(f_0=45\text{ Hz}\)。离线分析采用 1 s 滑动窗和 0.05 s 步长，并显式记录窗函数、相干增益与等效噪声带宽；这样，同频幅相估计与噪声分位数都能够追溯到确定的处理带宽。

设三轴磁力计某一轴的离散信号为

\[s(t) = A(t)\sin(\omega_0 t + \phi) + n(t)\]

其中 \(A(t)\) 为待提取的工频幅值包络、\(\phi\) 为相位、\(n(t)\) 为宽带噪声。构造两路正交参考 \(r_I(t) = \sin(\omega_0 t)\)、\(r_Q(t) = \cos(\omega_0 t)\)，分别与 \(s(t)\) 相乘并按积化和差展开：

\[\begin{aligned} x_I(t) &= s(t)\,r_I(t) = \frac{A(t)}{2}\left[\cos\phi - \cos(2\omega_0 t + \phi)\right] + n(t)\sin(\omega_0 t) \\ x_Q(t) &= s(t)\,r_Q(t) = \frac{A(t)}{2}\left[\sin\phi + \sin(2\omega_0 t + \phi)\right] + n(t)\cos(\omega_0 t) \end{aligned}\]

再经截止频率 \(f_c = 0.5\text{ Hz}\) 的四阶 Butterworth 低通滤波器滤除 \(2\omega_0\) 倍频分量与带外噪声，得到两路直流正交分量：

\[I(t) = \text{LPF}\{x_I(t)\} = \frac{A(t)}{2}\cos\phi, \qquad Q(t) = \text{LPF}\{x_Q(t)\} = \frac{A(t)}{2}\sin\phi\]

由此同时恢复幅值包络与相位：

\[R(t) = 2\sqrt{I(t)^2 + Q(t)^2} = A(t), \qquad \phi(t) = \arctan2\big(Q(t),\, I(t)\big)\]

DLIA 的三路输出恰好成为后续三个环节的枢纽：其一，三轴幅值包络合成为 \(|B_{f_0,\mathrm{rms}}(t)| = \sqrt{R_x^2 + R_y^2 + R_z^2}\)，直接供第 3.3.3 节的半高全宽深度反演使用；其二，相位角 \(\phi(t)\) 反映铠装涡流引起的相位滞后 \(\Delta\phi\)（对应第 2.3.3 节的屏蔽相位机制）；其三，由幅值与噪声底得到的瞬时信噪比 \(\text{SNR} = 20\log_{10}\dfrac{R(t)}{\sigma_{noise}}\) 直接驱动 EKF 的观测噪声协方差，SNR 高时 \(R_k\) 接近预设下界 \(R_{\min}\)，SNR 骤降时 \(R_k\) 增大、滤波器逐步转向 DVL 与惯性递推。这样，``预处理''与``融合''通过 SNR 这一标量真正连成一条闭环链路。

45 Hz 短背景记录的统一重分析进一步显示，矩形窗下向量幅值 p99 为 2.301 nT，改用 Hann 窗并线性去趋势后降至 0.398 nT，降幅约 82.7\%，说明有限记录的频谱泄漏是原始高噪声分位数的主要来源。对电缆与 ArUco 联合记录采用同一相干参考后，1049 个锁相窗口全部超过上述短背景阈值；其中 777 个位姿对齐点进入三轴复 I/Q 几何拟合，自由复尺度和固定峰值电流模型分别得到 \(R^2=0.8792\) 与 \(R^2=0.8711\)。这些结果把 DLIA 从接口设计推进为已有实测支撑的信号处理环节，并说明连续 I/Q、窗函数和参考频率必须作为磁观测的一部分共同记录。短背景时长较短所带来的统计限制集中在第 5.6 节讨论。

\section{基于改进 EKF 的声磁协同状态估计}\label{sec:ch03-3-3}

\subsection{系统状态向量与运动学模型定义}\label{sec:ch03-3-3-1}

在解决了``对得上、看得清''两层预处理问题之后，本节进入算法主体，状态向量、预测方程和观测方程的具体形式。理解这一节的关键是先回答一个看似简单但常被忽略的问题：状态向量到底应该包含哪些量？太少会导致系统模型不可观，太多则会引入冗余和数值病态。

本系统的状态向量按``运动状态 + 传感器误差''两类组织。运动状态包含 AUV 在 NED 坐标系下的位置 \texttt{p}（3 维）、单位四元数姿态 \texttt{q}（4 个标量、3 个自由度）和导航坐标系速度 \texttt{v}（3 维）；传感器误差包含 IMU 加速度计零偏 \texttt{b\_a}（3 维）和陀螺仪零偏 \texttt{b\_g}（3 维）。位置和姿态用于回答``AUV 在哪里、朝向哪里''，速度用于联系连续时间内的运动学方程，零偏则把``传感器永远存在缓变误差''这一物理事实纳入估计目标，只有这样，长时间航行积分误差才不会无限发散。

误差状态滤波中的标准状态分裂与误差注入过程可参见相关理论文献\cite{sola2017quaternion}。依照这一形式，本文把上述状态分为名义状态 \texttt{x\_nom}（直接维护数值）和误差状态 \texttt{δx}（始终在零附近波动）两部分。每次观测更新只对误差状态做线性化和卡尔曼增益计算，更新结束后再把误差状态``注入''名义状态并重置为零。四元数表示负责避免欧拉角奇异，误差状态结构则负责把姿态线性化限制在局部小扰动空间。

预测步骤在 IMU 输入下递推，名义状态按连续时间 IMU 积分公式更新；误差状态按一阶线性化的状态转移矩阵 \texttt{F} 演化，过程协方差 \texttt{Q} 累积。观测步骤在每个观测到达时调用，DVL 给出体坐标系下三轴速度修正、深度计给出 z 轴位置修正、磁场反演给出相对电缆的横向偏移和埋深线索、声呐特征给出电缆中心线方向。每个观测都有自己的雅可比矩阵 \texttt{H} 和观测协方差 \texttt{R}，后者会被自适应模块在线调整。

把上述描述写成形式化定义。名义状态由 16 个标量参数表示，但受单位四元数约束而具有 15 个自由度；各分量按位置、速度、姿态、加速度计零偏和陀螺仪零偏五组组织：

\[\boldsymbol{x} = \begin{bmatrix} \boldsymbol{p}^n & \boldsymbol{v}^n & \boldsymbol{q}_b^n & \boldsymbol{b}_a & \boldsymbol{b}_g \end{bmatrix}^T\]

其中 \(\boldsymbol{p}^n, \boldsymbol{v}^n \in \mathbb{R}^3\) 为 NED 系位置与速度，\(\boldsymbol{q}_b^n\) 为体系到导航系的姿态四元数，\(\boldsymbol{b}_a, \boldsymbol{b}_g \in \mathbb{R}^3\) 为加速度计与陀螺仪零偏。对应的 15 维误差状态在四元数处用三维旋转矢量 \(\delta\boldsymbol{\theta}\) 表示，从而回避姿态奇异：

\[\delta\boldsymbol{x} = \begin{bmatrix} \delta\boldsymbol{p} & \delta\boldsymbol{v} & \delta\boldsymbol{\theta} & \delta\boldsymbol{b}_a & \delta\boldsymbol{b}_g \end{bmatrix}^T \in \mathbb{R}^{15}\]

预测步以 IMU 比力 \(\boldsymbol{f}_b\) 与角速度 \(\boldsymbol{\omega}_b\) 递推名义状态（\(\boldsymbol{R}_b^n\) 为四元数对应的旋转矩阵、\(\boldsymbol{g}^n\) 为重力）：

\[\dot{\boldsymbol{p}}^n = \boldsymbol{v}^n, \qquad \dot{\boldsymbol{v}}^n = \boldsymbol{R}_b^n(\boldsymbol{f}_b - \boldsymbol{b}_a) + \boldsymbol{g}^n, \qquad \dot{\boldsymbol{q}}_b^n = \tfrac{1}{2}\,\boldsymbol{q}_b^n \otimes \begin{bmatrix} 0 \\ \boldsymbol{\omega}_b - \boldsymbol{b}_g \end{bmatrix}\]

误差状态按一阶线性化传播，其连续时间转移矩阵 \(\boldsymbol{F}\) 与过程噪声协方差 \(\boldsymbol{Q}\) 决定误差协方差 \(\boldsymbol{P}\) 的演化：

\[\delta\dot{\boldsymbol{x}} = \boldsymbol{F}\,\delta\boldsymbol{x} + \boldsymbol{G}\,\boldsymbol{w}, \qquad \dot{\boldsymbol{P}} = \boldsymbol{F}\boldsymbol{P} + \boldsymbol{P}\boldsymbol{F}^T + \boldsymbol{G}\boldsymbol{Q}\boldsymbol{G}^T\]

其中转移矩阵的非平凡分块来自速度对姿态误差、零偏的依赖：

\[\boldsymbol{F} = \begin{bmatrix} \boldsymbol{0} & \boldsymbol{I}_3 & \boldsymbol{0} & \boldsymbol{0} & \boldsymbol{0} \\ \boldsymbol{0} & \boldsymbol{0} & -\boldsymbol{R}_b^n[\boldsymbol{f}_b - \boldsymbol{b}_a]_\times & -\boldsymbol{R}_b^n & \boldsymbol{0} \\ \boldsymbol{0} & \boldsymbol{0} & -[\boldsymbol{\omega}_b - \boldsymbol{b}_g]_\times & \boldsymbol{0} & -\boldsymbol{I}_3 \\ \boldsymbol{0} & \boldsymbol{0} & \boldsymbol{0} & \boldsymbol{0} & \boldsymbol{0} \\ \boldsymbol{0} & \boldsymbol{0} & \boldsymbol{0} & \boldsymbol{0} & \boldsymbol{0} \end{bmatrix}\]

式中 \([\cdot]_\times\) 为反对称叉乘矩阵。观测更新采用统一形式：给定观测 \(\boldsymbol{z}_k\) 与观测模型 \(h(\boldsymbol{x})\)，新息（innovation）、新息协方差、卡尔曼增益与误差状态修正依次为

\[\boldsymbol{y}_k = \boldsymbol{z}_k - h(\hat{\boldsymbol{x}}_k^-), \quad \boldsymbol{S}_k = \boldsymbol{H}_k\boldsymbol{P}_k^-\boldsymbol{H}_k^T + \boldsymbol{R}_k, \quad \boldsymbol{K}_k = \boldsymbol{P}_k^-\boldsymbol{H}_k^T\boldsymbol{S}_k^{-1}, \quad \delta\hat{\boldsymbol{x}} = \boldsymbol{K}_k\boldsymbol{y}_k\]

更新完成后把 \(\delta\hat{\boldsymbol{x}}\) 注入名义状态（位置/速度/零偏直接相加、姿态按 \(\boldsymbol{q} \leftarrow \boldsymbol{q}\otimes\delta\boldsymbol{q}(\delta\boldsymbol{\theta})\) 复合）并将误差状态重置为零。各观测通道的雅可比 \(\boldsymbol{H}_k\) 分别为：DVL 观测体系速度，\(\boldsymbol{H}_{dvl} = \partial(\boldsymbol{R}_n^b\boldsymbol{v}^n)/\partial\delta\boldsymbol{x}\)；深度计观测 \(z\) 位置，\(\boldsymbol{H}_{dep} = [0,0,1,\,\boldsymbol{0}_{1\times12}]\)；声学与磁学观测模型的雅可比在 §3.3.2、§3.3.3 分别给出。磁学通道中，峰值时刻提供横向中心线约束，FWHM 提供垂直距离或埋深约束，二者共同抑制 AUV 相对电缆的位置漂移。

\subsection{声学观测模型：基于侧扫声呐特征的管线定位}\label{sec:ch03-3-3-2}

为什么把声学观测放在第一个？因为它直接给出几何信息，电缆在 AUV 视野中的位置和方向，这是 ES-EKF 横向位置估计最直接的修正源。当声呐图像中电缆边缘清晰时，可以从图像中提取电缆中心线，进而计算 AUV 相对于电缆的横向偏移 \texttt{d\_lat} 和相对方位角 \texttt{θ\_rel}。

声学观测模型的输入是从声呐图像中提取的特征点序列（电缆中心线在 AUV 体坐标系下的位置），输出是观测向量 \texttt{z\_sonar\ =\ {[}d\_lat,\ θ\_rel{]}\^{}T}。观测方程把这两个量与状态向量中的位置和姿态联系起来，如果 AUV 已知自己的全局位置和姿态，又已知电缆的全局路由，那么 \texttt{d\_lat} 和 \texttt{θ\_rel} 是状态的确定性函数。雅可比矩阵 \texttt{H\_sonar} 由这一函数对状态求偏导得到。

但是，声呐观测的成立依赖一个前提：图像质量足够支持电缆边缘的稳定提取。这个前提在以下情况下会被打破，浑浊水域中悬浮颗粒造成的散射、底质突变区域的强背景反射、电缆局部掩埋或被沙泥覆盖、声呐多径反射在窄水道中造成虚像。一旦这些情况出现，特征提取算法或者输出错误的``假阳性边缘''，或者完全输出空集合。如果状态估计器仍然按``声呐总是可信''处理，错误观测会被强行注入，造成估计跳变甚至发散。

针对这种情况，本系统在声呐管线的输出端附加一个``特征质量分数''，它综合了边缘强度、连续性、与历史预测的一致性等指标，输出一个 {[}0,1{]} 区间的标量。这个分数被映射为观测协方差 \texttt{R\_sonar} 的缩放因子：质量分数高时 \texttt{R\_sonar} 接近名义值，质量分数低时 \texttt{R\_sonar} 被放大数十倍甚至完全屏蔽该观测。这样，``何时使用声呐''不再是一个二值开关，而是一个连续可调的权重，整个滤波器的行为也更平滑。

\subsection{磁学观测模型：基于毕奥-萨伐尔定律的埋深反演}\label{sec:ch03-3-3-3}

磁学观测在物理上与声学完全不同，它不是看几何，而是看电流。海缆中流动的工频电流在周围空间产生一个由毕奥-萨伐尔定律描述的磁场\cite{griffiths2017electrodynamics,liu2024submarine}，这个磁场强度与电流幅值、距离的二次方倒数和导线相对方位角有关。如果 AUV 上磁传感器测得的三轴磁场分量记为 \texttt{B\ =\ {[}B\_x,\ B\_y,\ B\_z{]}\^{}T}，那么它就是电缆位置参数（横向偏移 \texttt{d\_lat}、埋深 \texttt{d\_depth}、相对走向 \texttt{θ\_cable}）的非线性函数。

直接三轴矢量反演可通过最小化观测残差 \texttt{\textbar{}\textbar{}B\_obs\ -\ B\_predicted(d\_lat,\ d\_depth,\ θ\_cable)\textbar{}\textbar{}²} 求解电缆位置参数，但该基线在单时刻下存在多解，对电流幅值、铠装衰减和螺旋相位也较敏感。连续轨迹虽然能增加约束，仍不能消除幅值尺度不确定性。因此，本文不把直接矢量滑窗反演作为主输入链，而把它保留为物理模型对照。

但物理上的难点在于本章 3.1 节已经反复强调的几件事：三相抵消使得观测幅值仅为单芯电缆的几分之一甚至更弱、铠装屏蔽进一步通过涡流衰减磁场、螺旋绞合使磁场幅值随路径距离呈周期性而非单调变化。这三个效应叠加意味着``在电缆正上方''和``略偏离电缆''的磁场差别可能不大，反演的几何敏感度变弱。本文据此采用两个工程化处理：第一，在反演前先做``空间相关性匹配滤波''，利用 AUV 运动轨迹与电缆走向的几何先验，把磁场时间序列与``假设电缆位于某处时的预期序列''做相关计算；第二，把磁观测的雅可比矩阵在三相相位附近做``光滑化''处理，避免梯度在抵消区附近出现尖锐的零点。

实际主链路依次为``三轴模拟量采集与标定 → DLIA 幅相提取 → 旋转不变模量包络 → FWHM 拟合 → 横向/垂直观测与质量度量 → ES-EKF 更新''。该链路不直接使用三轴矢量方向，而采用第 2.3.4 节定义的模量包络 \(|B(t)| = \sqrt{R_x^2 + R_y^2 + R_z^2}\)，并调用第 2.3.5 节的半高全宽关系把它转化为几何观测量。回顾该模型：当 AUV 以航速 \(v\)、跨切角 \(\alpha\) 之字形横切电缆时，模量包络呈洛伦兹型

\[|B(x)| = \frac{\mu_0 m_{eq}}{2\pi\,(x^2 + z^2)}, \qquad x(t) = v\sin\alpha\,(t - t_0)\]

在电缆正上方取得峰值，其半高全宽满足 \(x_{half} = \pm z\)，因此电缆垂直距离由半高时间跨度 \(\Delta t_{FWHM}\) 直接给出：

\[z_{mag} = \frac{1}{2}\,v|\sin\alpha|\cdot\Delta t_{\mathrm{FWHM}}, \qquad d_{depth} = z_{mag} - h_{alt}\]

其中 \(h_{alt}\) 为高度计读数。这一反演的关键优势在于它\textbf{完全解耦了对未知电流 \(I\)、铠装衰减率 \(\eta\) 与磁偶极矩 \(m_{eq}\) 的依赖}，这些正是海况下最难先验标定的量。据此，磁学观测向 ES-EKF 提供两个绝对约束：横向中心线位置（峰值时刻 \(x = 0\)）与垂直距离 \(z_{mag}\)。前者修正横向偏移 \(d_{lat}\)，后者作为埋深观测 \(z_{mag}\) 并入 §3.3.1 的观测向量。相应的观测方程与雅可比为

\[\boldsymbol{z}_{mag} = h_{mag}(\boldsymbol{x}) + \boldsymbol{v}_{mag}, \qquad \boldsymbol{H}_{mag} = \frac{\partial h_{mag}}{\partial \delta\boldsymbol{x}}\bigg|_{\hat{\boldsymbol{x}}_k^-}\]

在第 2.3.5 节的模型假设成立时，\(z_{mag}\) 是不依赖幅值尺度参数的几何观测量；其观测噪声协方差 \(R_{mag}\) 主要由 DLIA 输出的瞬时 SNR、横切速度与交角误差以及 FWHM 拟合残差决定。峰值位置与 \(z_{mag}\) 分别约束横向中心线和垂直距离，不能把垂直距离单独解释为对 DVL 横向漂移的直接校正。

第 5.4 节的 45 Hz 电缆--TMR--ArUco 联合实验从幅值与相位两方面验证了毕奥--萨伐尔几何映射；FWHM 深度反演则依赖横切速度、交角和幅值包络的共同可辨识性，其系统级精度应通过可溯源电流与距离基准评定。

上述从原始三轴采样到 ES-EKF 观测的完整链路可归纳为图~\ref{fig:ch03-cable-perception-pipeline}。图中左列为在线磁链：三轴采样经标定去趋势后进入 DLIA，输出旋转不变模量与 SNR，再由峰值定位与 FWHM 得到横偏、垂直距离、埋深及质量分；声呐则以质量分驱动的稀疏几何锚点接入电缆几何级在线修正，与载体级 ES-EKF 位姿共同解算横偏、走向与埋深，最终汇入部署侧 \texttt{CableTrackingOutput} 供行为树与 UA-MPC 消费。真值仅进入离线评价，不参与在线管线。该结构对应第 2.3.5 节主动横切构造几何可观测性的原理（图~\ref{fig:ch03-active-crossing}）：一次横切同时给出横向中心线与垂直距离两类观测。

\begin{figure}[htbp]
\centering
\includegraphics[width=0.92\linewidth,height=\textheight,keepaspectratio,alt={声磁协同电缆感知流水线：三轴采样经DLIA与FWHM得到几何观测，声呐提供稀疏锚点，两级估计汇入部署侧输出}]{architecture/cable_acoustic_magnetic_perception_pipeline.png}
\caption{声磁协同电缆感知流水线：原始三轴磁场经 DLIA 与半高全宽转化为几何观测，声呐提供质量分驱动的稀疏锚点，载体级 ES-EKF 与电缆几何级在线修正两级估计汇入部署侧 \texttt{CableTrackingOutput}}\label{fig:ch03-cable-perception-pipeline}
\end{figure}

\subsection{声磁接力与动态权重分配逻辑}\label{sec:ch03-3-3-4}

到这里为止，声学和磁学两个观测通道已经各自定义清楚。已有水下多传感器融合研究表明，异构观测的可靠性差异与失效切换是系统设计必须处理的基础问题\cite{shaukat2021multisensor}。在此基础上，本文需要进一步回答：当两者都可信时如何加权？当两者都不可信时如何处理？当一个可信一个不可信时如何切换？本节从这三种情形出发，说明本文声磁接力与动态权重机制的具体设计。

第一种情形是``声呐主导，磁场辅助''。在大部分浅埋或部分裸露的电缆段，声呐能提供清晰的几何信息（横向偏移 \texttt{d\_lat} 精度 0.2 m 量级），磁场反演的精度则较粗糙（因为磁场对几何的敏感度有限）。此时声呐观测的 \texttt{R\_sonar} 较小、磁场观测的 \texttt{R\_mag} 较大，滤波器主要按声呐修正横向位置，磁场只起辅助校验作用，例如检验``是否真的有电流流过''，避免被海底其他线状目标（管道残骸、绳索）误识别。

第二种情形是``磁场主导，声呐沉默''。当电缆深埋（\textgreater1 m）或被沙泥完全覆盖时，声呐几乎无法分辨电缆，输出端的特征质量分数接近零，滤波器自动把 \texttt{R\_sonar} 放大到几乎屏蔽。此时磁场观测成为唯一的电缆相关线索，虽然几何精度不如声呐，但至少能给出``AUV 是否在电缆走向附近''的二值判断和粗略的横向偏移估计。

第三种情形是``两者都不可信''。在 DVL 同时丢包、磁畸变同时发生的极端工况下，声磁两个通道都被自适应模块降权，系统只能退回到 IMU 递推 + 深度修正的基础模式。此时滤波器的协方差会持续增长，触发置信度衰减机制，进而由行为树切换到搜索模式或紧急上浮。这种``协方差驱动的退化''让系统在感知能力不足时不会强行作业，符合工程上的安全优先原则。

把这三种情形组合起来，``接力''指的是观测主导权随环境变化而平滑切换，``动态权重''指的是切换过程通过协方差 \texttt{R} 的连续调整完成，而不是硬阈值开关。这种设计避免了硬切换可能引入的突变性误差，使滤波器输出在状态空间内始终连续。

要把上述定性逻辑落到工程可实现的判据上，还需要给出两类通道的失效边界与量化触发条件。声呐侧的失效边界由其物理机理决定：侧扫声呐（540/850 kHz）在电缆裸露或位于开阔沟槽时成像清晰，但一旦掩埋深度超过约 10 cm，其对电缆的识别率骤降至接近 0\%，浅水区还会因多径产生伪影；浅地层剖面仪（4--24 kHz）虽能穿透泥下 0.5--3 m，但只有在 AUV 近乎正上方飞过时才准确，横向漂移超过约 1 m 即立即失锁，且无法用于大面积搜索。磁场侧则存在另一类几何模糊：在急弯段，单磁力计仅能观测标量幅值 \(|B|\)，无法据此判断电缆是左拐、右拐还是埋深突然变大，因而在弯道段需要主动横向机动来重建方向可辨识性。需要说明的是，这里的``弯道''是相对物理机理的定性描述，而非固定失效阈值；专用算法仓库对纯磁段做定半径弧扫描（详见第 5.5.11 节）时，直到 30 m 环境硬下限之前均未观测到弯道触发的失效，说明只要在线先验修正与主动横切在位，纯磁在弯道的瓶颈更接近电缆几何本身的物理下限，而非某一固定曲率半径。

据此，本系统把接力过程组织为一个由自适应协方差 \(\boldsymbol{R}_k = \text{diag}(R_{mag}, R_{sonar})\) 驱动的决策流程，覆盖三类典型工况：工况 A（裸露、水质好）声呐质量分数 \(S_{sonar} \approx 1.0\)，取 \(R_{sonar}\) 小、\(R_{mag}\) 中，以声呐为主实现米级贴线跟踪；工况 B（泥下掩埋）\(S_{sonar} \to 0\)、\(R_{sonar} \to \infty\)、\(R_{mag}\) 收小，主导权接力给磁力计；工况 C（急弯段方向模糊）状态协方差 \(\boldsymbol{P}_k\) 因方向不可辨识而增大，触发行为树执行小角度之字形横切，用第 2.3.5 节的 FWHM 定理重新解算方向与深度。这一流程把``何时机动''从随机搜索变成由协方差判据触发的确定性动作。

在此基础上可以刻画系统的两条可靠性边界，二者均由专用算法仓库的单次纯仿真扫描给出，是同源算法的边界摘要而非主仓端到端或实测结论。其一是稀疏声呐鲁棒性：即便声呐有效率降到约 20\%（等效丢包 80\%、刷新率降至 0.4 Hz），稀疏锚点配合磁路径与在线先验修正仍可维持长程任务推进，在急弯段依靠磁场包络加 DVL 递推撑过声呐盲区，跟踪段横偏维持在米级量级而非厘米级。其二是声呐初始锁定后中断的纯磁续航：当声呐仅用于初始锁定、进入跟踪后被强制离线时，受控之字形与在线先验修正可在所测工况中维持推进，跟踪段横偏约 1.7--2.6 m。但必须同时指出其负边界：当声呐从一开始就完全缺席（无初始锁定）时，磁模量与局部过线证据不足以单独完成全局切向与方向初始化，系统无法完成初始捕获，因此``全盲重捕''并不成立，这一工况应作为可观测性负边界保留。上述边界共同界定了第 5 章接力消融实验的三个对照组：声呐 100\%（理想）、声呐 20\%（主推，验证严重丢包下仍可跟踪）、声呐初始锁定后中断（退化，验证纯磁直道可跑、弯道触发行为树主动横切）。

\section{感知不确定性的量化度量模型}\label{sec:ch03-3-4}

\subsection{基于信噪比的观测协方差 R 在线调整}\label{sec:ch03-3-4-1}

到目前为止，本章已经数次提到``自适应 R''，但这个 R 究竟应该如何在线调整？这一节给出具体的量化模型。整套思路可以归纳为一句话：用归一化新息平方（normalized innovation squared, NIS）做``事后体检''，用观测质量分数做``事前预判''，二者协同决定每个观测通道的协方差缩放系数。

先看 NIS 的物理含义\cite{barshalom2001estimation}。在卡尔曼滤波框架下，每次观测都会产生新息 \texttt{y\ =\allowbreak{}\ z\ -\ H·x\_\allowbreak{}pred}，并伴随理论新息协方差 \texttt{S\ =\ H·P·H\^{}T\ +\ R}。标准 NIS 定义为 \texttt{y\^{}T·S\^{}(-1)·y}；若观测维数为 \(m\) 且模型自洽，其期望为 \(m\)。为比较不同维数的观测通道，本文另定义维数归一化 NIS 为标准 NIS 除以 \(m\)，其理论期望为 1。二者必须分开报告：前者用于卡方门控，后者用于跨通道或跨场景汇总。

但是，单次 NIS 的随机波动较大，持续偏离或窗口统计更适合判断协方差是否长期一致；仅依赖它处理单帧质量突降容易产生滞后。这就是为什么本系统在 NIS 之外还引入``事前预判''的观测质量分数。声呐质量分数来自图像特征提取的边缘强度和连续性，磁场质量分数来自滤波后的瞬时 SNR，DVL 质量分数则来自底部跟踪有效性标志和速度残差。每个分数都是一个 {[}0, 1{]} 区间的标量，再通过预设的非线性映射（例如指数函数）转换为 R 缩放因子 \(\alpha \in [1,\alpha_{\max}]\)，最终的在线观测协方差为 \texttt{R\_\allowbreak{}online\ =\allowbreak{}\ α\ ·\ R\_\allowbreak{}nominal}。

把``事前预判''和``事后体检''组合起来：每次观测到达时，先按质量分数 \texttt{α} 调整 R，给滤波器做一次``先验加权''；累积一段时间后，再按观测源分别用 NIS 做后验复核。持续高于对应自由度的卡方上界表示协方差偏小，持续低于下界则表示协方差过于保守。分源处理使 DVL 三维观测和深度一维观测拥有各自的统计参照，避免不同自由度的原始 NIS 在同一平均值中相互抵消。

把上述闭环写成形式化判据。对每次观测计算新息向量与新息协方差，再分别构造标准 NIS 与维数归一化 NIS：

\[\boldsymbol{y}_k = \boldsymbol{z}_k - \boldsymbol{H}_k\hat{\boldsymbol{x}}_k^-, \qquad \boldsymbol{S}_k = \boldsymbol{H}_k\boldsymbol{P}_k^-\boldsymbol{H}_k^T + \boldsymbol{R}_k,\]

\[\epsilon_k=\boldsymbol{y}_k^T\boldsymbol{S}_k^{-1}\boldsymbol{y}_k,\qquad \bar{\epsilon}_k=\frac{\epsilon_k}{m}.\]

在模型自洽的假设下 \(\epsilon_k\sim\chi^2(m)\)，据此取显著性水平 \(\alpha=0.05\) 的上侧门限 \(\epsilon_{\mathrm{th}}=\chi^2_{1-\alpha}(m)\)。当 \(\epsilon_k\) 超过门限时判定该观测为可疑残差，并按残差超标幅度对观测协方差做指数膨胀（软打折），从而平滑压低卡尔曼增益\cite{piche2016online,chen2023kalman}：

\[\boldsymbol{R}_{k,adaptive} = \frac{\boldsymbol{R}_k}{c_k}, \qquad c_k = \exp\!\big(-\lambda\,\max(0,\ \epsilon_k-\epsilon_{\mathrm{th}})\big)\]

其中 \(\lambda > 0\) 为衰减速率。残差正常（\(\epsilon_k\le\epsilon_{\mathrm{th}}\)）时 \(c_k=1\)、协方差不变；残差严重超标时 \(c_k\to0\)、\(\boldsymbol{R}_{k,adaptive}\) 增大、卡尔曼增益降低，即该观测被渐进降权而非硬性丢弃。为避免少量脉冲残差在窗口统计尚未收敛时直接支配状态修正，本文还实现 Huber 型即时降权。令 \(r_k=\sqrt{\epsilon_k/m}\)，当 \(r_k>\delta\) 时按 \(w_k=\delta/r_k\) 放大当前观测协方差为 \(\boldsymbol{R}_{k,robust}=\boldsymbol{R}_{k,adaptive}/w_k\)，否则 \(w_k=1\)。窗口级分源缩放处理持续失配，Huber 权重处理单帧异常，两者的时间尺度和统计角色互不替代。

这一``分源卡方诊断 + 窗口软打折 + 即时鲁棒权重''构成 §3.3.4 接力逻辑在协方差层面的实现。对 24 个运行包的 35243 次真实量测更新审计表明，深度与 DVL 的一致性偏差方向相反，因而历史实现中混合两种自由度的固定阈值只能解释为操作性触发，不能解释为统一卡方检验。第 5.5.5 节进一步用固定 R、历史全局缩放、分源缩放和 Huber 更新作同输入对照，并把时间戳噪声膨胀与真正的延迟状态处理区分开来。

\subsection{传感器失效检测与置信度衰减机制}\label{sec:ch03-3-4-2}

3.4.1 节解决了``观测协方差怎么调''，但还有一个更上层的问题：当一个传感器整体失效（不是观测质量稍降，而是连续多个周期都不可用）时，系统应该如何把这一信息传递给行为树和控制层\cite{sun2018variational}？这就是置信度衰减机制要解决的问题。

机制的核心是把每个传感器通道的协方差状态映射为一个标量置信度 \(c \in [0,1]\)。具体而言，本系统通过 \texttt{cov\_to\_conf} 函数完成这一映射，它把当前观测协方差与名义协方差的比值取对数后做 sigmoid 平滑，使得：当协方差接近名义值时 \texttt{c\ ≈\ 1.0}，当协方差被放大到名义值的 10 倍时 \texttt{c\ ≈\ 0.5}，当协方差超过名义值 100 倍或观测连续多个周期不可用时 \texttt{c\ →\ 0}。

为什么用 sigmoid 而不是线性或者硬阈值？答案是工程性的，sigmoid 的两端饱和、中间陡峭的特性正好符合``在大部分时间稳定、在过渡区敏感、在极端情况下饱和''的工程需求。如果用线性映射，极端工况下置信度会跌到负值或者过快归零；如果用硬阈值，过渡区会出现``在 0.6 和 1.0 之间反复跳变''的颤振。sigmoid 在数学上是这两种极端的折中。

置信度衰减的输出有两个去向。第一个去向是 UA-MPC，置信度低时，跟踪误差权重通过非线性函数被放大，控制努力代价则通过 sigmoid 平滑降低，为必要纠偏保留控制权威；指令平滑性由速率和带宽约束保证。第二个去向是行为树，置信度低于第一阈值时降速，低于第二阈值时切换到搜索模式，低于第三阈值时触发紧急上浮。三档阈值依次是 0.6、0.3、0.1，对应``小心、警觉、退出''三种响应等级。

整体上看，置信度衰减把``传感器是否可信''这一原本只在滤波器内部使用的信息，转化为一个被多个上层模块共享的接口标量。这种设计既避免了上层逻辑直接读取协方差矩阵带来的耦合复杂度，也使``不确定性感知''贯穿整个系统而不只是停留在状态估计这一层。

\section{仿真环境下的状态估计验证}\label{sec:ch03-3-5}

\subsection{针对 DVL 丢包的鲁棒性测试}\label{sec:ch03-3-5-1}

理论说完，下一步考察多种传感器扰动下的估计连续性。

DVL 丢包是水下任务中最常见的失效场景，也是验证 ES-EKF 鲁棒性最直接的对象。本系统通过 {\small\path{scenario_dvl_dropout_{10,30,60,90}.yaml}} 配置四档丢包率（10\%、30\%、60\%、90\%），每档场景中 ChaosInjector 会在指定时间窗口内冻结 DVL 输出，模拟浊水或杂草遮挡。这种设计的好处是丢包时机可重复、丢包概率精确可控，便于做多 seed 统计。

{\small\path{docs/experiment/benchmark_test_log.md}} 中已有 baseline 单次（n=1）的定位对比结果，三种状态估计方法（ES-EKF、Std-EKF、Raw DR）在干净环境下的精度排序符合理论预期：ES-EKF 优于 Std-EKF，Std-EKF 优于 Raw DR。这一结果支持``工具链可用、ES-EKF 在 baseline 下有效''的初步结论，但不能直接外推到丢包场景下的鲁棒性。

更进一步，第 5 章中已完成的``P1 sensor 3 seed sweep''（24 个 run）覆盖了 DVL dropout 四档 × 3 seed，结果在第 5.5.5 节中给出。从那里的数据可以看到：四档 dropout 下 ES-EKF 均达到 100\% 成功率（24/24 ok），XY RMSE 在 2.76--3.55 m 区间内，并不严格随丢包率单调上升。这意味着 30 s 片段、seed 初始化和观测时序对结果影响显著，丢包率与定位精度之间不是简单的``扰动越强误差越大''关系。

四档丢包结果共同说明，当前自适应估计器能够在观测间歇失效时维持状态输出；误差对丢包率的非单调响应则表明，短时窗口中的观测时序同样是决定结果的重要因素。为把``随机种子效应''与``更新时序效应''分开，第 5.5.5 节进一步采用相同噪声序列，对 0 ms 与 100 ms 两种 DVL 更新相位作五种子成对比较。两档保持相同的 5 Hz 更新频率，仅改变观测落在预测周期中的位置；300 s 结果显示，相位变化本身即可产生与场景扰动同量级的误差差异。这一现象从状态可观性角度可以解释：DVL 约束体坐标速度、深度计约束垂向位置，但长期水下航段缺少持续的绝对水平位置与航向观测，预测误差在不同更新相位下被投影到速度和姿态误差子空间的方式不同，最终通过积分放大为水平漂移。因此，多速率融合不能只声明各传感器频率，还必须把时间戳、相位和插值策略作为估计器定义的一部分。

\subsection{磁信号畸变下的定位精度分析}\label{sec:ch03-3-5-2}

磁畸变场景通过 {\small\path{scenario_mag_distortion_{light,heavy}.yaml}} 配置，磁饱和阈值从 1e-6 T 收紧到 1e-8 T。该场景测试磁场观测不可靠时，状态估计是否能通过自适应模块降低磁观测权重，并维持估计连续性。其物理意图是模拟两类常见情形：一是 AUV 靠近母船等大磁源时的近场干扰，二是磁传感器自身量程受限导致的硬饱和。

第 5 章 P1 sweep 中包含磁畸变（heavy）单场景的 3 seed 结果，XY RMSE 约 3.42\allowbreak{}±0.10 m，与 baseline（约 4.62 m）的对比显示磁畸变情况下 ES-EKF 反而略好，但这并非``磁畸变改善定位''的反直觉结论，而是 30 s 片段下声学和 DVL 观测的主导作用使磁畸变影响被掩盖。本文据此把磁畸变场景的核心价值定位为``验证滤波器在磁观测降权时不会发散''，而不是``量化磁畸变对定位精度的具体影响''。

轻、重两档磁畸变以及综合压力场景均已纳入三种子结果。该组实验验证的是磁观测降权时的估计连续性；真实频谱、相位畸变和软饱和将由实测噪声回放进一步补充，统一见第 5.6 节。

\section{本章小结}\label{sec:ch03-3-6}

本章按``问题定义 → 框架设计 → 算法主体 → 不确定性量化 → 仿真验证''的顺序，建立了水下电缆巡检任务下的声磁协同状态估计链路。多速率时间对齐和杆臂/安装角补偿保证观测在统一时空基准下进入滤波器；45 Hz 相敏解调从三轴磁场中提取连续 I/Q、幅值和 SNR；误差状态扩展卡尔曼滤波再以分源观测协方差连接声呐、磁场、DVL 与惯性递推。

八类扰动场景、三种子共二十四次运行均保持稳定输出，说明该架构能够在 DVL 丢包、磁畸变、声呐杂波及联合扰动下维持估计连续性。分源 NIS 审计进一步揭示深度观测噪声低估与 DVL 观测过度保守同时存在；五分钟成对实验又确认观测相位是独立于丢包比例的重要变量。由此，不确定性接口从单纯的触发机制推进到可按观测源和时间结构诊断的统计量，本章形成了从实测磁信号预处理、几何观测构造到协方差驱动融合的完整方法链；其集中适用范围与后续验证见第 5.6 节。
